\section{Local distance, completion and the metric tangent cone}
\label{sec:metric-cone}

Normalize the arrows by $T_p(A,B)=(\sqrt{c_+}A,\sqrt{c_-}B)$.
The balanced set is
\[
 \mathcal Z=\{(X,Y):XX^*=Y^*Y,\ YY^*=X^*X\},
\]
with real metric $2(\|\dd X\|_{\rm HS}^2+\|\dd Y\|_{\rm HS}^2)$.
Let $\mathsf W=(\ZZ/2)^2\rtimes S_2$ act on $\CC^2$
by independent signs and permutation. For
$u=(u_1,u_2)$ and $U=\diag(u_1,u_2)$, define
\begin{equation}\label{eq:actual-moduli-parametrization}
 \mathfrak F_p([u])=
       [\mathcal A_p(\chi_p(T_p^{-1}(U,U)))].
\end{equation}
The brackets on the right denote actual unitary gauge
classes, equivalently actual holomorphic bundle classes
at these HE points.
Choose one $R_0>0$ so that $|u|\le R_0$ lies strictly
inside every construction domain, uniformly in $p$.
Put $\mathcal M_p(R)=\mathfrak F_p(\overline B_R/\mathsf W)$
for $0<R\le R_0$.

An admissible path in this locus is continuous and has
a finite subdivision into smooth paths in its quotient
orbit-type strata, with improper endpoint lengths allowed.
Its length is the integral of the actual harmonic
connection norm of its bundle-deformation tangent.
Let $d_{p,R}$ be the infimum of these lengths among
paths confined to $\mathcal M_p(R)$.

\begin{theorem}\label{thm:metric-cone}
The map \eqref{eq:actual-moduli-parametrization} is a
homeomorphism of $\overline B_R/\mathsf W$ onto its
actual image, with smooth image-stratum parametrizations.
On every stratum its harmonic metric satisfies
\begin{equation}\label{eq:diagonal-harmonic-metric}
       (1-C|u|)\,4|\dd u|^2
       \le\mathfrak F_p^*g_{\rm harm}\le(1+C|u|)\,4|\dd u|^2.
\end{equation}
For $|u|,|v|\le R$, set
$d_{\rm flat}([u],[v])=2\min_{\gamma\in\mathsf W}|u-\gamma v|$.
Then
\begin{equation}\label{eq:confined-distance-comparison}
 \sqrt{1-CR}\,d_{\rm flat}([u],[v])
 \le d_{p,R}(\mathfrak F_p[u],\mathfrak F_p[v])
 \le\sqrt{1+CR}\,d_{\rm flat}([u],[v]).
\end{equation}
The open-radius locus completes isometrically to this
closed-radius locus. On the unit quotient ball,
\begin{equation}\label{eq:rescaled-distance}
 \sup_{[u],[v]}
 \left|\rho^{-1}d_{p,\rho}(\mathfrak F_p[\rho u],\mathfrak F_p[\rho v])
                          -d_{\rm flat}([u],[v])\right|\le C\rho.
\end{equation}
The same estimate holds with the single fixed-locus
distance $d_{p,R_0}$ for $\rho<R_0/8$.
The actual local metric space therefore has the unique
pointed metric tangent cone
\begin{equation}\label{eq:moduli-cone}
                         (\CC^2,4|\dd u|^2)/\mathsf W
\end{equation}
at the split bundle. These comparisons and the resulting
Gromov--Hausdorff convergence are uniform over the smooth
conifold base, including simultaneous approaches to its
discriminant. The distances use paths in the specified
local locus.
\end{theorem}

\subsection{The exact finite quotient}

The balanced normal form in Section~\ref{sec:balancing}
is $X=\diag(d_j)$, $Y=\diag(d_je^{\mathrm i\theta_j})$.
Choose $u_j=d_je^{\mathrm i\theta_j/2}$.
Acting by $k_+=\diag(e^{\mathrm i\theta_j/2})$, $k_-=I$
makes both arrows equal $U$. The phase at a zero entry
is immaterial. Under the compact action $XY$ changes by
conjugation. For $(U,U)$ its eigenvalue multiset is
$\{u_1^2,u_2^2\}$. Thus two section points can be
equivalent only if they differ by $\mathsf W$.
Conversely permutations are realized simultaneously in
both blocks, and each sign change by a diagonal sign
in the plus block alone. This proves the quotient
identification, including repeated and zero entries.
On a closed radius ball it is a homeomorphism by
compactness, since
$\|X\|_{\rm HS}^2+\|Y\|_{\rm HS}^2=2|u|^2$ and
the compact group quotient is Hausdorff.

The four strata are distinct nonzero squares, repeated
nonzero squares, one zero square, and the origin.
Their preimages are the complement of the four complex
reflecting hyperplanes, the nonzero parts of $u_1=\pm u_2$,
the nonzero parts of the coordinate axes, and the origin.
They agree with the actual factor and stabilizer types.

The section also gives smooth local coordinates on each
stratum. On the balanced set $XY$ is normal and
$XX^*=|XY|$, since
$(XY)(XY)^*=(XX^*)^2=(XY)^*(XY)$.
For distinct eigenvalues its spectral projections vary
smoothly. For a nonzero eigenvalue $\lambda$, choose a
smooth local square root $u$ and a smooth unit eigenvector
$e_+$; then $e_-=X^*e_+/\bar u$ is unit and satisfies
$Xe_-=ue_+$, $Ye_+=ue_-$.
At a simple zero eigenvalue use smooth frames of the
kernels, on which both arrows vanish. On the repeated
nonzero stratum $XY=\lambda I$ and the same formula
works with an arbitrary smooth plus frame. These are
qualitative coordinates within a stratum; no uniform
bound for their derivatives near another stratum is used.

For a diagonal section tangent $(V,V)$ the ambient norm
is $4\|V\|_{\rm HS}^2$. Its pairing with a vertical is
\[
 2\operatorname{Re}\tr
       V^*(\xi_+U-U\xi_-+\xi_-U-U\xi_+)
   =2\operatorname{Re}\tr V^*[\xi_++\xi_-,U]=0.
\]
Thus these tangents are horizontal, and the stratum
coordinates just proved exhaust the quotient tangents.
The exact finite quotient metric is $4\sum_j|\dd u_j|^2$.
On the repeated stratum it is $8|\dd u|^2$, and on
the one-nonzero stratum it is $4|\dd u|^2$; the origin
has zero stratum tangent.

Its path distance on a closed ball is $d_{\rm flat}$.
Local lifts of stratum paths and the triangle inequality
give the lower bound. For the upper bound choose a
minimizing $\gamma$ and project the straight segment
from $u$ to $\gamma v$. Convexity keeps it in the ball.
Each reflecting hyperplane meets the segment at at most
one parameter unless it contains the segment; a finite
subdivision therefore makes the projected segment
admissible. Its length attains the stated lower bound.

\subsection{Actual harmonic lengths and completion}

Theorem~\ref{thm:unitary-rigidity} and the compact quotient
argument of Section~\ref{sec:harmonic-quotient} identify
the finite quotient homeomorphically with actual classes.
Each image stratum is represented by a smooth family
of integrable HE connections. Its harmonic tangent
has no kernel beyond the compact vertical directions
by Theorem~\ref{thm:harmonic-quotient}.
Since $|T_p^{-1}(U,U)|\le C|u|$, that theorem and the
exact horizontal metric prove
\eqref{eq:diagonal-harmonic-metric}.
Integrating it over any admissible path gives the lower
bound in \eqref{eq:confined-distance-comparison}; using
the projected straight segment gives the upper bound.
In particular this is a finite nondegenerate distance
with exactly the compact quotient topology.

For interior endpoints the open-locus path distance
equals the restriction of the closed-locus distance.
One inequality is immediate. For the other, contract
any finite-length closed-locus path by $u\mapsto(1-\epsilon)u$.
The contracted pieces remain in their strata. The actual
tensor is smooth there, so their speeds converge pointwise.
Equation~\eqref{eq:diagonal-harmonic-metric} dominates
them by a constant times the original Euclidean speed,
which is integrable by the lower bound in that equation.
Dominated convergence proves convergence of lengths,
including improper endpoint integrals. Add short radial
paths at the two interior endpoints; their lengths tend
to zero. Infimizing gives equality of the distances.
Radial contraction also proves density of the open
locus in the complete compact closed locus, establishing
the claimed completion.

Homogeneity of $d_{\rm flat}$ and
\eqref{eq:confined-distance-comparison} with $R=\rho$
prove \eqref{eq:rescaled-distance}. The flat unit quotient
ball has diameter at most four. The identity correspondence
therefore has distortion $O(\rho)$, giving the uniform
Gromov--Hausdorff bound directly.

\subsection{Localization in a fixed local metric space}

Shrink $R_0$ so that $CR_0<1/2$.
For endpoints of parameter norm at most $\rho$, the
straight-path upper bound is less than $5\rho$.
The radius $|u|$ is well defined on the quotient, and
its variation is at most half the finite length.
Hence a path starting at radius at most $\rho$ and
reaching $8\rho$ has actual length at least
\[
             2\sqrt{1-CR_0}(8\rho-\rho)>9\rho.
\]
Every path of length at most $6\rho$ between these
endpoints stays inside radius $8\rho$. Near-minimizing
paths can be chosen with this length bound. Along them
the lower tensor factor is $1-8C\rho$, while the
straight segment supplies the upper factor $1+C\rho$.
Thus the fixed-locus distance satisfies the same relative
comparison with error $O(\rho)$ for these endpoints.
This proves its version of \eqref{eq:rescaled-distance}.

To pass from parameter balls to actual metric balls,
first the coarse comparison puts every small metric
ball inside a comparable parameter ball. The sharper
fixed-locus estimate then gives
\[
          d_{p,R_0}(0,\mathfrak F_p[u])=2|u|(1+O(|u|)).
\]
An actual ball of radius $2\rho$ lies between parameter
balls of radii $\rho(1-C'\rho)$ and $\rho(1+C'\rho)$.
Radial projection between these changes distances by
$O(\rho^2)$. After rescaling, the distortion is $O(\rho)$.
The same argument on every fixed multiple of $\rho$
gives compatible pointed convergence to the full cone
\eqref{eq:moduli-cone}, proving uniqueness of the tangent.

Equivalently this cone is the symmetric square of
$(\CC,4|\dd u|^2)/\{\pm1\}$, whose cone angle is $\pi$.
The invariant coordinates $\lambda_j=u_j^2$ do not
carry an ordinary Euclidean metric at zero.
The theorem includes all local polystable factor types
and every simultaneous admissible base approach. It
concerns the harmonic metric on this actual local locus;
it does not compare shortcuts through a larger moduli
space or assert a metric on a singular base fibre.
